% Artículo VI. Propietario: indice_articulos. Adición autónoma de exposición.
% Procedencia: RESULTADO_RECUPERADO; controles documentales nuevos separados.
% REV2/fuente/modulos/15_comparacion_cuantitativa_integral.md (completo).
% REV2/fuente/modulos/16_precision_robustez_y_contraste.md (completo).
% REV2/fuente/modulos/19_realizacion_exhaustiva_estado_por_estado.md (completo).
% Integral: integracion_83/masas_p0/corpus/source/masas/
% 19_comparacion_cuantitativa_posterior.tex, parte anterior a \endinput.
% Integral: integracion_83/masas_p0/tex/apendice_inventario_471_384_sucesor.tex.
% Depende de la construcción de atlas, firma, torres y realización material
% incorporada por el editor en las secciones precedentes. No la sustituye.

\section{Réalisation des observables et comparaison métrologique
}
\label{sec:vi-contraste-metrologico}

La valeur d’une masse acquiert son sens au sein de la structure qui la
produit et de la représentation dans laquelle elle est mesurée. La construction
APP--TRIT--TPK détermine d’abord les trajectoires admissibles, leurs registres
orientés et les lecteurs de la signature ; les tours et la réalisation matérielle
agissent sur ces objets. La comparaison expérimentale se situe au terme de
cette composition. Sa fonction consiste à comparer des observables du même
type, à conserver l’identité de l’état et à quantifier l’accord dans une
coordinatisation commune. Cette position permet d’étudier conjointement la structure
mathématique, la coordonnée évaluée et son interprétation physique.


La présente section réunit trois niveaux qui remplissent des fonctions différentes :
le domaine interne des routes et de leurs quotients ; les classes et les évaluations
individualisées dans les propriétaires des masses ; et l’inventaire extérieur des
observables utilisé comme référence. L’ordre d’exposition conserve la
direction constructive de l’article. Les grandeurs extérieures apparaissent comme
termes de comparaison, tandis que les valeurs internes conservent les signatures,
les opérations et les conditions qui fixent leur évaluation.


\subsection{Domaine structurel, classes matérielles et inventaire extérieur
}
\label{subsec:vi-dominios-contraste}

Soient \(\mathcal C\) l’ensemble des cellules APP et \(\mathcal R\) leur domaine
de routes orientées. L’identification d’une route contient les coordonnées
de cellule et les deux orientations initiales :

\[
 \gamma=(r,c,h_0,v_0),\qquad
 r,c\in\{1,\ldots,9\},\quad h_0,v_0\in\{-1,+1\}.
\]
Par conséquent, \(\#\mathcal C=81\) et \(\#\mathcal R=81\cdot4=324\).
Les \(56\) collections de trajectoires et les \(13\) multisections appartiennent
à des quotients et à des regroupements structurels du même domaine, avec les
relations et les lecteurs établis précédemment. Ces quatre cardinaux
décrivent des objets internes distincts. L’attribution d’une représentation
physique peut conserver une fibre complète ; une étiquette d’espèce
n’impose pas à elle seule un choix ponctuel au sein de celle-ci.


Le tableau explicite de \(23\) classes organise les réalisations par secteur.
Sa fonction est de relier la résidence interne à la représentation, à la
charge, au spin, à la couleur, à la saveur et au type d’observable. Une classe peut
comprendre plusieurs espèces, et une identité catalographique peut réunir
différents états de charge. La correspondance s’exprime au moyen d’applications
typées et conserve ces multiplicités.


L’inventaire métrologique figé contient \(471\) enregistrements de masse et
\(384\) de largeur. Le corpus documente \(1170\) identités dans \(438\)
groupes, dont \(403\) contiennent au moins une de ces observables.
La somme

\[
 471+384=855
\]
compte les observables extérieures de tableau récapitulatif. L’atlas \(81/56/13/324\),
les \(23\) classes et les \(855\) enregistrements appartiennent ainsi à trois niveaux
de description. Leur association conserve le type de chaque ensemble, au lieu
de transformer leurs cardinaux en un nombre unique de particules ou de prédictions.


\begin{center}
\begin{tabular}{p{0.49\linewidth}rrr}
\toprule
Collection catalographique
 & Groupes
 & Masses
 & Largeurs
\\
\midrule
Général
 & 46 & 53 & 9\\
Quarks
 & 6 & 8 & 1\\
Mésons
 & 201 & 243 & 225\\
Baryons
 & 149 & 166 & 149\\
Enregistrement gravitationnel extérieur
 & 1 & 1 & 0\\
\midrule
Total avec observable
 & 403 & 471 & 384\\
\bottomrule
\end{tabular}
\end{center}

La ligne gravitationnelle conserve l’identité qui figure dans la source extérieure ;
sa présence dans l’inventaire ne définit pas une particule élémentaire de HMT.
Les \(35\) groupes restants conservent leur identité catalographique, même si
l’état arrêté ne leur attribue pas de masse ou de largeur de tableau récapitulatif. L’intégrité
documentaire englobe les deux faits : les grandeurs consignées et les identités
pour lesquelles cet état arrêté n’exprime pas l’une de ces grandeurs.


\subsection{L’observable avant la coordonnée numérique
}
\label{subsec:vi-tipos-observable}

Une sortie matérielle se spécifie au moyen de l’état, de la représentation,
de l’opérateur et du procédé de lecture. Pour une fibre \(F_X\), soit
\(\widehat M_X\) l’opérateur de masse déjà construit et soit \(P_{X,a}\)
le projecteur correspondant à l’état ou au canal de préparation \(a\).
Pour \(\Psi_X\) normalisé dans son image, la distribution spectrale est

\[
 \mu_{X,a}(B)
 =\langle\Psi_X,
   P_{X,a}E_{\widehat M_X}(B)P_{X,a}\Psi_X\rangle,
\]
où \(E_{\widehat M_X}\) est la mesure spectrale de l’opérateur. Cette
expression conserve l’état préparé et la multiplicité de la fibre.
Une coordonnée scalaire unique pour tout état préparé dans ce sous-espace
s’obtient sous la condition

\[
 \operatorname{Ran}(P_{X,a})\subset\operatorname{Dom}(\widehat M_X),
 \qquad \widehat M_XP_{X,a}=m_XP_{X,a}.
\]
En effet, chaque vecteur de ce sous-espace est alors un vecteur propre de
l’opérateur de masse pour la même valeur \(m_X\), et sa mesure spectrale est
concentrée en cette valeur. Si l’on emploie l’égalité de compression
\(P_{X,a}\widehat M_XP_{X,a}=m_XP_{X,a}\), on exige en outre que
\(P_{X,a}\) réduise \(\widehat M_X\) ; en dimension finie, cela équivaut
à leur commutation. Sans cette condition, la compression peut fixer uniquement
la valeur moyenne. Lorsque la mesure spectrale maintient plusieurs valeurs, la
sortie se présente comme une distribution ou un multiplet avec ses poids.


Pour une densité fixe \(\rho\) dans une fibre finie, soient \(m_j\)
les valeurs propres distinctes et \(E_j\) leurs projecteurs spectraux.
Les poids \(p_j=\operatorname{tr}(\rho E_j)\geq0\), avec
\(\sum_jp_j=1\), déterminent

\[
 \overline m=\sum_jp_jm_j,\qquad
 \operatorname{Var}_{\mu}(m)=\sum_jp_j(m_j-\overline m)^2.
\]
\paragraph{Critère de concentration scalaire pour l’état fixé.
}
La variance s’annule si et seulement si toutes les valeurs de poids positif
coïncident. En effet, chaque terme est positif ou nul ; la somme est nulle
précisément lorsque \(m_j=\overline m\) pour tout \(j\) tel que
\(p_j>0\). On détermine ainsi quand une distribution de masse admet
une représentation scalaire sans perte de la dispersion de l’état préparé.
Ce critère concerne la densité fixée et se distingue de la condition
opératorielle sur tout le sous-espace. Le cas électronique utilise son
opérateur sur la fibre centrale, tel qu’il a été construit précédemment ;
le rang un d’un projecteur arbitraire ne suffit pas à concentrer la mesure
spectrale d’un autre opérateur.


Masse et largeur utilisent des lecteurs coordonnés de types différents. Pour un
état résonant, une convention de pôle peut s’exprimer dans les coordonnées
d’énergie sous la forme

\[
 z_X=M_Xc_{\rm int}^{2}-\frac{i}{2}\Gamma_X.
\]
La partie réelle est la coordonnée de masse-énergie ; \(\Gamma_X\) a la
dimension d’une énergie et mesure la largeur. Une fois fixée la coordinatisation temporelle,
\(\tau_X=\hbar_{\rm int}/\Gamma_X\). La relation suppose le régime de
pôle et de survie pour lequel elle a été construite ; un taux, une largeur et
une durée de vie maintiennent leurs unités et leurs règles de conversion. Cette
distinction explique pourquoi l’inventaire conserve deux classes d’observables.


Le même principe détermine la comparaison par secteur. Les masses courantes
de quark s’expriment avec le schéma \(\mathscr S\) et l’échelle \(\mu\) ;
les résonances distinguent masse de pôle, paramétrage de Breit--Wigner et
convention sur couche. Les neutrinos distinguent valeurs propres de masse,
différences de carrés, saveurs et paramètres de mélange. Les bornes
unilatérales conservent le niveau de confiance et la convention statistique
documentés. Une borne supérieure ne devient pas une valeur centrale positive,
ni une différence de carrés une masse absolue.


\subsection{Fixation de la construction et correspondance des registres
}
\label{subsec:vi-congelacion}

La réalisation d’une espèce \(X\) transmet l’information nécessaire
à partir de ses antécédents internes. Nous pouvons la réunir, à des fins de comparaison,
dans le registre

\[
 \mathfrak D_X=
 \bigl(F_X,\rho_X,P_{X,a},\mathcal L_X,
       \sigma_X,\eta_X,\widehat M_X,\mathcal U_X\bigr),
\]
où \(\rho_X\) identifie la représentation, \(\mathcal L_X\) le
registre généalogique, \(\sigma_X\in\mathbb Z^5\) la signature,
\(\eta_X\) la section de tour et \(\mathcal U_X\) la coordinatisation des
unités. Toutes ces données reçoivent leur sens des opérations
APP--TRIT--TPK et des réalisations démontrées dans l’article. Les réunir
dans un n-uplet est une opération de documentation : leurs constructions et
leurs preuves restent dans les sections qui les produisent.


La référence extérieure est représentée séparément par

\[
 o=(\mathrm{id},\mathrm{espèce},\mathrm{observable},
      v,\mathrm{unité},\mathscr S,\mu,
      \mathrm{incertitude},\mathrm{limite},\mathrm{confiance},
      \mathrm{source},\mathrm{édition}).
\]
L’association \((\mathfrak D_X,o)\) doit s’établir au moyen de
l’identité physique et du type d’observable. Dans un tableau étendu, cette
opération se réalise par clés d’identité et de provenance ; la proximité
de deux chiffres ne remplit pas la fonction de clé. Le registre maintient
l’étiquette de particule ou d’antiparticule, le groupe qui possède l’observable et
l’identité de la propriété mesurée.


L’indépendance du procédé de sélection s’examine en masquant ou en
permutant la colonne métrologique et en vérifiant que \(\mathfrak D_X\)
reste inchangé. Ce test a un sens pour un générateur dont la
composition effective peut être exécutée ; la simple conservation d’un fichier
de sortie ne constitue pas à elle seule un tel test. L’empreinte cryptographique
fixe l’artefact comparé et rend les modifications détectables, mais
l’indépendance causale appartient à l’opération et à ses dépendances.


Les statuts sont attribués au domaine et à l’opération qui fondent chaque
affirmation. Une ligne peut documenter simultanément un opérateur interne,
une évaluation étant donnée une signature et une comparaison extérieure, sans identifier
ces trois niveaux entre eux. Toute réalisation ultérieure est liée à la
définition, à la proposition ou au tableau qui établit sa portée. En particulier, les
classifications documentaires du catalogue s’appliquent à l’état arrêté que spécifie
chaque enregistrement et ne modifient pas le statut mathématique d’une construction
ultérieure.


\subsection{Correspondance métrologique des vingt-trois classes
}
\label{subsec:vi-clases-metrologia}

Le tableau suivant conserve les \(23\) classes de la source et spécifie la
forme de lecture qu’exige chacune lorsqu’elle atteint l’interface métrologique. Sa
fonction diffère de la classification opératorielle de la
proposition~\ref{vi:prop:producto-fibrado-especie-ruta} et de l’inventaire
documentaire des tableaux~\ref{tab:vi-estados-clases} et~\ref{tab:vi-clases} :
ici sont distinguées la grandeur comparable, la convention de réalisation et
les données qui doivent l’accompagner. Les dénominations de saveur et de secteur
orientent la correspondance ; l’objet mathématique transmis reste dans
les fibres, les représentations et les lecteurs correspondants.


\begingroup
\small
\setlength{\tabcolsep}{3pt}
\begin{table}[!htbp]
\centering\fontsize{9.2}{10.6}\selectfont\renewcommand{\arraystretch}{1.04}
\caption{Lecture métrologique des vingt-trois classes et type d’objet transmis.
}
\label{tab:vi-clases-metrologia}
\begin{tabular}{p{0.20\linewidth}p{0.25\linewidth}p{0.47\linewidth}}
\toprule
Classe
 & Représentants
 & Objet documenté pour la lecture
\\
\midrule

Lepton chargé 1
 & Électron et positron
 & Route centrale, conjugaison et coordonnée bêta.
\\
Lepton chargé 2
 & Muon et antimuon
 & Fibre, lecteurs et évaluation de la signature déclarée.
\\
Lepton chargé 3
 & Tau et antitau
 & Fibre, lecteurs et évaluation de la signature déclarée.
\\
Neutrino 1
 & Saveur électronique et conjuguée
 & Fibre neutre et réalisation masse--saveur.
\\
Neutrino 2
 & Saveur muonique et conjuguée
 & Fibre neutre et réalisation masse--saveur.
\\
Neutrino 3
 & Saveur tauique et conjuguée
 & Fibre neutre et réalisation masse--saveur.
\\
Quark arriba 1 & \(u,\bar u\) & Fibre de saveur et de couleur ; coordinatisation de schéma et d’échelle.
\\
Quark abajo 1 & \(d,\bar d\) & Fibre de saveur et de couleur ; coordinatisation de schéma et d’échelle.
\\
Quark arriba 2 & \(c,\bar c\) & Fibre de saveur et de couleur ; coordinatisation de schéma et d’échelle.
\\
Quark abajo 2 & \(s,\bar s\) & Fibre de saveur et de couleur ; coordinatisation de schéma et d’échelle.
\\
Quark arriba 3 & \(t,\bar t\) & Fibre de saveur et de couleur ; coordinatisation de schéma et d’échelle.
\\
Quark abajo 3 & \(b,\bar b\) & Fibre de saveur et de couleur ; coordinatisation de schéma et d’échelle.
\\
Jauge électromagnétique
 & Photon
 & Coordinatisation nulle et représentation d’hélicité.
\\
Jauge chromatique
 & Gluons
 & Coordinatisation nulle et représentation adjointe de couleur.
\\
Calibre cargado & \(W^+,W^-\) & Route, signature et convention de pôle de la réalisation.
\\
Calibre neutro & \(Z^0\) & Route, mélange neutre et convention de pôle.
\\
Scalaire
 & Higgs
 & Représentation scalaire, signature et convention de pôle.
\\
Composé mésonique
 & Mésons et conjugués
 & Constituants, registre de liaison et lecture composée.
\\
Composé baryonique
 & Baryons et conjugués
 & Composition trilinéaire, liaison et lecture composée.
\\
Topologique Fibonacci
 & \(1,\tau\) & Anneau de fusion et dimension de Frobenius--Perron.
\\
Topologique Ising
 & \(1,\psi,\sigma\) & Règles de fusion et représentation sectorielle.
\\
Topologique cyclique
 & \(\mathbb Z/6\mathbb Z\) & Caractères d’enroulement et représentation des tresses.
\\
Virtuel
 & Sections intermédiaires & Horizon fini, prolongement et obstruction.
\\
\bottomrule
\end{tabular}
\end{table}
\endgroup

Dans les leptons chargés, la lecture conduit à la masse de la section centrale
ou à l’évaluation de la fibre sélectionnée. Dans les neutrinos, elle conserve en outre
la transformation entre masse et saveur et les différences quadratiques ; c’est pourquoi
elle ne se réduit pas à un scalaire isolé. Dans les quarks, chaque chiffre est accompagné du
schéma et de l’échelle de renormalisation. Les classes électromagnétique et
chromatique expriment respectivement une borne expérimentale et la section nulle
de jauge ; les classes \(W\), \(Z\) et scalaire requièrent la convention de pôle
ou de résonance déclarée. Les composés conservent masse, largeur et identité
des constituants. Les trois classes topologiques sont comparées au moyen de leurs
lois de fusion, de tressage et d’écart une fois fixée la réalisation dynamique ;
leurs dimensions et leurs caractères appartiennent à des codomaines propres. La classe
virtuelle conserve le pôle, l’obstruction et l’horizon de prolongement. Cette
lecture détermine quel objet atteint la métrologie sans confondre les
vingt-trois classes avec vingt-trois scalaires universels de masse.


\subsection{État métrologique arrêté, unités et exactitude de la transcription
}
\label{subsec:vi-corte-metrologico}

On utilise l’état identifié par le corpus comme Particle Data Group
2026, avec une date d’arrêté au \(15\) janvier \(2026\) et une consultation
documentée du \(26\) juillet \(2026\). Cette section reproduit cet
état documentaire : elle n’incorpore pas de mise à jour ultérieure du catalogue.
Les références du supplément d’évaluation et celles de la grille historique
maintiennent leur provenance individuelle ; elles ne reçoivent pas automatiquement l’édition
de l’inventaire général.


Chaque valeur est conservée comme chaîne décimale originale, avec sa
présentation, son unité, son incertitude et sa clé de provenance. Cette décision
évite qu’une conversion intermédiaire en virgule flottante ne change les derniers
chiffres ou n’élimine des zéros significatifs de la transcription. La chaîne
originale et la représentation typographique sont deux données distinctes : la
première permet de reproduire exactement l’arithmétique du fichier ; la
seconde communique le résultat avec une précision adaptée à la lecture.
Les chiffres de représentation numérique d’une référence n’augmentent pas sa
précision expérimentale.


La conversion entre MeV et GeV utilise le facteur exact \(10^3\), appliqué
à la valeur et à son incertitude. Pour présenter une masse, l’unité s’écrit
MeV/\(c^2\) ou GeV/\(c^2\) ; pour présenter l’énergie de repos,
MeV ou GeV. La relation \(E=mc^2\) est la coordinatisation qui permet de passer d’une
lecture à l’autre. Une comparaison de coordonnées exige de fixer cette coordinatisation dans
les deux colonnes.


L’état arrêté de l’inventaire maintient expressément comme non fournis le
schéma et l’échelle de ses enregistrements extérieurs. Le nom d’une propriété
peut contenir une indication telle que ``Breit--Wigner'' ; elle est conservée littéralement,
mais cette indication partielle ne complète pas les métadonnées que la comparaison
requiert. Les cases vides, les marques de donnée non disponible et celles de
non-applicabilité restent différenciées. Une inconnue documentaire
n’est pas remplacée par zéro.


\subsection{Évaluations individualisées et conditions de lecture
}
\label{subsec:vi-evaluaciones-documentadas}

Le tableau~\ref{tab:vi-evaluaciones-veinte} réunit \(20\) présentations
numériques : la coordonnée bêta électronique, deux sorties nulles et
\(17\) évaluations conditionnées par des descripteurs déclarés. Parmi
ces dernières, \(11\) correspondent à des secteurs élémentaires et \(6\) à des
composés. Ce décompte exprime le contenu de cette présentation documentaire ;
les fibres, les collections et les registres du domaine conservent leur extension propre.


Le tableau maintient les vingt entrées. Pour faciliter sa lecture, on
identifie quatre modalités : \(\mathrm{B}\), coordonnée bêta
électronique ; \(\mathrm{N}\), coordinatisation nulle ; \(\mathrm{F}\), évaluation
conditionnée par une signature complète ; \(\mathrm{A}\), évaluation angulaire
conditionnée. L’indice \(c\) signale la composition matérielle. La
condition de chaque entrée appartient à son type de résultat et est conservée
également dans le catalogue électronique qui accompagne le texte.


\begingroup
\small
\setlength{\tabcolsep}{3pt}
\begin{table}[!htbp]
\centering\fontsize{8.8}{10.2}\selectfont
\setlength{\tabcolsep}{4pt}\renewcommand{\arraystretch}{1.04}
\caption{Vingt évaluations individualisées, avec leur coordinatisation et leur modalité de comparaison.
}
\label{tab:vi-evaluaciones-veinte}
\begin{tabular}{p{0.13\linewidth}p{0.32\linewidth}p{0.38\linewidth}p{0.07\linewidth}}

\toprule
État
 & Coordonnée HMT, MeV/\(c^2\)
 & Référence documentaire ultérieure
 & Type
\\
\midrule

\(e^\pm\) & \(0.510998950\allowbreak690000809\) & \(0.51099895069(16)\) MeV/\(c^2\) & B\\
\(\mu^\pm\) & \(105.658360294\allowbreak435829303\) & \(105.6583755\pm0.0000023\) MeV/\(c^2\) & F\\
\(\tau^\pm\) & \(1776.860917631\allowbreak432295595\) & \(1776.93\pm0.09\) MeV/\(c^2\) & F\\
\(u,\bar u\) & \(2.165924536\allowbreak173907663\) & \(2.16\pm0.07\) MeV/\(c^2\) & A\\
\(d,\bar d\) & \(4.679550162\allowbreak948874172\) & \(4.70\pm0.07\) MeV/\(c^2\) & A\\
\(c,\bar c\) & \(1270.500838641\allowbreak911135286\) & \(1.2729\pm0.0045\) GeV/\(c^2\) & A\\
\(s,\bar s\) & \(92.940093907\allowbreak845640641\) & \(92.9\pm0.7\) MeV/\(c^2\) & A\\
\(t,\bar t\) & \(172597.479544019\allowbreak136261367\) & \(172.60\pm0.27\) GeV/\(c^2\) & A\\
\(b,\bar b\) & \(4190.119763299\allowbreak790218075\) & \(4.186\pm0.006\) GeV/\(c^2\) & A\\
\(\gamma\) & \(0\) & \(<10^{-18}\) eV/\(c^2\), borne conservée
 & N\\
\(g\) & \(0\) & Coordinatisation de jauge nulle ; sans ligne expérimentale de comparaison
 & N\\
\(W^\pm\) & \(80378.964909704\allowbreak547024865\) & \(80.3625\pm0.0077\) GeV/\(c^2\) & F\\
\(Z^0\) & \(91187.533444313\allowbreak407844855\) & \(91.1879\pm0.0020\) GeV/\(c^2\) & F\\
\(H\) & \(125292.466042536\allowbreak133000433\) & \(125.13\pm0.11\) GeV/\(c^2\) & A\\
\(\pi^0\) & \(134.976724327\allowbreak872125739\) & \(134.976827767685\) MeV/\(c^2\), supplément
 & \(\mathrm F_c\)\\
\(\pi^\pm\) & \(139.570485171\allowbreak572191375\) & \(139.570390983681\) MeV/\(c^2\), supplément
 & \(\mathrm F_c\)\\
\(K^\pm\) & \(493.677085649\allowbreak530612476\) & \(493.676599458041\) MeV/\(c^2\), supplément
 & \(\mathrm F_c\)\\
\(K^0\) & \(497.611186497\allowbreak750457359\) & \(497.610767531258\) MeV/\(c^2\), supplément
 & \(\mathrm F_c\)\\
\(p,\bar p\) & \(938.271751546\allowbreak984898299\) & \(938.27208943\) MeV/\(c^2\), supplément
 & \(\mathrm F_c\)\\
\(n,\bar n\) & \(939.565810472\allowbreak507023313\) & \(939.56542194\) MeV/\(c^2\), supplément
 & \(\mathrm F_c\)\\
\end{tabular}
\end{table}
\endgroup

Dans ce tableau, les valeurs de référence sont présentées avec la notation
décimale et les unités déclarées pour chaque observable. Le registre de
calcul conserve en outre la chaîne décimale employée dans l’évaluation,
y compris la différence entre cette chaîne et sa présentation arrondie pour
le tau, \(Z\) et Higgs. La comparaison des quarks
reste informative tant que la coordinatisation de schéma et d’échelle n’est pas réunie.
Les secteurs \(W,Z,H\) requièrent de conserver également la convention de
résonance. Les six références composées proviennent du supplément
identifié dans les registres d’évaluation ; leurs métadonnées ne sont pas complétées
par analogie avec les références élémentaires.


Les signatures complètes documentées dans les modalités \(\mathrm F\) et
\(\mathrm F_c\), dans l’ordre\linebreak
\(\sigma=(n_A,s_C,k_{\Delta_4},b_{120},b_\zeta)\), sont :

\[
\begin{array}{c|r@{,\,}r@{,\,}r@{,\,}r@{,\,}r}
\mu&42&-3&8&0&2\\
\tau&64&0&25&-1&6\\
W&94&-1&0&-2&4\\
Z&95&-1&-11&0&-4\\
\pi^0&44&-4&-25&1&-2\\
\pi^\pm&44&1&-2&2&-1\\
K^\pm&54&-1&17&-2&2\\
K^0&54&0&32&3&-2\\
p&59&0&9&1&0\\
n&59&1&-34&0&1
\end{array}
\]
Les modalités angulaires conservent, respectivement pour
\(u,d,c,s,t,b,H\), les paires

\[
(n_A,s_C)=(11,7),(17,8),(61,8),(41,-3),(100,-1),(71,-5),(97,9).
\]
Les deux composantes déclarées ne sont pas étendues avec
trois zéros implicites. L’équation d’évaluation et le mécanisme qui
produit un descripteur sont des affirmations différentes ; les preuves des deux
se trouvent dans leurs définitions et propositions respectives.


\subsection{Différences relatives et propagation de l’incertitude
}
\label{subsec:vi-residuales}

Une coordinatisation commune et deux coordonnées positives de la même observable étant fixées,
la différence relative et la différence logarithmique sont

\[
 \delta_X=\frac{m_X^{\rm HMT}-m_X^{\rm ext}}{m_X^{\rm ext}},
 \qquad r_X=\log\frac{m_X^{\rm HMT}}{m_X^{\rm ext}},
 \qquad \delta_X^{\rm ppm}=10^6\delta_X.
\]
Une conversion commune d’unités \(m\mapsto am\), avec \(a>0\),
laisse les deux différences invariantes : le facteur \(a\) s’annule dans les
quotients. La différence absolue est, elle, multipliée par \(a\), de même que
son incertitude. Pour une valeur extérieure nulle, une borne unilatérale ou un
intervalle sans estimation centrale, le quotient précédent est remplacé par
la comparaison correspondant à ce type d’observation.


Lorsqu’on dispose d’un modèle conjoint d’incertitude, soit
\(d_X=m_X^{\rm HMT}-m_X^{\rm ext}\). Sa variance est

\[
 u^2(d_X)=u^2(m_X^{\rm HMT})+u^2(m_X^{\rm ext})
          -2\operatorname{Cov}(m_X^{\rm HMT},m_X^{\rm ext}).
\]
La formule résulte du développement de
\(\operatorname{Var}(U-V)\). Si \(u(d_X)>0\), la différence
normalisée est \(z_X=d_X/u(d_X)\). Son interprétation statistique exige
le modèle probabiliste et la covariance correspondants ; un champ non
renseigné n’équivaut pas à une variance nulle. Les incertitudes asymétriques
sont conservées comme telles et ne sont pas automatiquement transformées en un écart
symétrique.


Pour plusieurs espèces, la matrice de covariance conjointe est maintenue. Les
signatures discrètes étant fixées, si \(f_X(\vartheta)=\log(m_X/m_{\rm ref})\)
dépend de coordonnées continues déjà construites et si \(J\) est sa matrice
jacobienne, la propagation au premier ordre donne

\[
 \Sigma_f=J\Sigma_\vartheta J^{\mathsf T}.
\]
Cette formule conserve les corrélations que partagent la paire angulaire,
\(\Delta_4\), le rapport d’action et les tours. Son application présuppose
le régime dans lequel la linéarisation contrôle les termes restants.
Un changement entier de signature est enregistré comme variation discrète de la
construction, non comme une fluctuation gaussienne infinitésimale.


La précision d’une évaluation a donc plusieurs sources : la
référence expérimentale ; la coordinatisation des unités, du schéma et de l’échelle ; la
distribution des routes ; l’approximation de la tour ; et le couplage
sélectionné. Si \(\Sigma_{ij}\) désigne la covariance entre les
contributions \(i,j\), la covariance totale de leur somme est
\(\sum_{i,j}\Sigma_{ij}\). Maintenir les blocs croisés évite
d’attribuer une indépendance à plusieurs coïncidences qui utilisent les mêmes
antécédents.


\subsection{Échelle absolue et contrôle de l’approximation harmonique
}
\label{subsec:vi-precision-torres}

L’échelle absolue et la correction relative interviennent à des endroits
différents de la construction. Une section basale de référence
\(m_\star^{(0)}\), construite en interne, étant fixée, soit
\(\chi_\star=\chi_{\rm full}(\sigma_\star,\eta_\star)\).
Dans la normalisation relative à cette section, l’expression est

\[
 m_X^\beta=m_\star^{(0)}\,
 \frac{\chi_{\rm full}(\sigma_X,\eta_X)}{\chi_\star}
 R_{\rm act}^{\kappa_X}.
\]
Le facteur basal commun reste explicite. Si la référence est électronique,
\(m_\star^{(0)}=m_e^{(0)}\) est la réalisation dimensionnelle de
\(\mathfrak e_0\), avec tous les facteurs de son équation basale, et
\(\chi_\star=\chi_{\rm full}(\sigma_e,\eta_e)\). Ainsi, on récupère
\(m_e^\beta=m_e^{(0)}R_{\rm act}^{\kappa_e}\). Si l’on utilise l’horloge
commune pour écrire \(m_\star^{(0)}=m_{\rm clk}b_\star\), le
facteur \(b_\star\) conserve cette normalisation basale ; on ne le remplace pas
par \(1\) et on ne l’incorpore pas deux fois au caractère. Pour deux états
dans la même coordinatisation et avec la même référence, le quotient est

\[
 \frac{m_Y^\beta}{m_X^\beta}
 =\frac{\chi_{\rm full}(\sigma_Y,\eta_Y)}
        {\chi_{\rm full}(\sigma_X,\eta_X)}
     R_{\rm act}^{\kappa_Y-\kappa_X}.
\]
L’annulation de \(m_\star^{(0)}/\chi_\star\) est exacte. Elle comprend
le préfacteur basal commun et les facteurs communs de la branche absolue
d’action, d’horloge et de coordinatisation de vitesse. Par
conséquent, la décade d’action \(10^{-34}\) conserve son rôle dans
la réalisation dimensionnelle et n’est pas insérée de nouveau dans le correcteur
relatif. Les différences entre routes restent dans leurs signatures, lecteurs
et sections de tour.


Pour contrôler une troncature, il convient de fixer la forme normale du caractère.
Soient

\[
 \zeta_{270}=135\Delta_4^2,\qquad
 N_h=\eta_h+b_{120}\mathbf 1_{\{h=120\}},\qquad
 \mathcal S=\langle90,120\rangle\setminus\{0\}.
\]
Alors
\[
\begin{split}
 \log\chi_{\rm full}
 ={}&n_AA_{\rm rad}+\frac{s_C}{6}C^*_{\rm rad}
       +k_{\Delta_4}\Delta_4+b_\zeta\zeta_{270}
       +\sum_{h\in\mathcal S}N_hs_h .
\end{split}
\]
Dans cette convention, le terme de hauteur \(120\) n’apparaît qu’une seule fois.
Le terme \(b_\zeta\zeta_{270}\) maintient son origine quadratique dans le
registre ; \(N_{270}s_{270}\) appartient à la lecture harmonique.
L’égalité de l’indice numérique n’identifie pas les deux objets.


\paragraph{Borne de troncature et erreur multiplicative.
}
Supposons que, pour \(h>H\), les coefficients déjà déterminés d’une
route satisfassent \(|N_hs_h|\leq Cq^h\), avec \(C\geq0\) et
\(0<q<1\). La queue logarithmique \(R_H\) vérifie

\[
 |R_H|\leq\sum_{\substack{h\in\mathcal S\\h>H}}Cq^h
 \leq\sum_{h=H+1}^{\infty}Cq^h
 =\frac{Cq^{H+1}}{1-q}=:\varepsilon_H.
\]
Si \(m_H\) est la coordonnée obtenue en tronquant cette tour et si les autres
facteurs restent fixes, \(m/m_H=e^{R_H}\). La monotonie de
l’exponentielle produit

\[
 e^{-\varepsilon_H}\leq\frac{m}{m_H}\leq e^{\varepsilon_H},
 \qquad
 \left|\frac{m-m_H}{m_H}\right|\leq e^{\varepsilon_H}-1.
\]
La borne permet d’attribuer une précision à l’évaluation à partir de la route et de ses
opérateurs. Les valeurs de \(C,q,H\) sont justifiées au moyen de cette
construction ; le résidu expérimental ne sélectionne ni la longueur de la
tour ni la borne que l’on entend démontrer.


\subsection{Conservation intégrale et reproductibilité de la comparaison
}
\label{subsec:vi-catalogo-reproducible}

Le catalogue conserve les enregistrements complets dans des formats tabulaires et
structurés. Chaque ligne maintient le fichier de provenance, son ordinal,
la clé originale et le numéro d’occurrence de cette clé. Cette
identification conserve même deux lignes littéralement identiques :
l’égalité de leurs contenus n’élimine pas la seconde occurrence. Le changement
de format préserve les chaînes originales, les cases vides et les
états, et joint l’empreinte des octets d’entrée.


L’inventaire de \(855\) observables est en outre relié aux lignes
métrologiques du corpus intégral. La correspondance est vérifiée
par la paire formée du type d’observable et de son identifiant ; le
contrôle compare les multiplicités, et pas seulement les ensembles de clés. Les tableaux
internes de cellules, de collections, de multisections et de routes restent séparés
de cet inventaire. Les \(23\) classes, les \(20\) présentations
numériques et la grille historique conservent également leurs fichiers et leurs statuts
propres.


Les contrôles de transcription vérifient la récupération de tous les
champs, l’ordre des lignes, les doublons, les chaînes décimales et la
correspondance avec les tableaux de référence. Leur portée est documentaire.
Les preuves mathématiques de l’atlas, des lecteurs et de la réalisation matérielle
sont celles qui établissent les sorties internes ; la comparaison métrologique
complète cet exposé au moyen d’observables et de coordinatisations bien définies.
L’intégrité documentaire et la suffisance d’une démonstration sont
vérifiées séparément.


L’extension du catalogue suit donc une règle cumulative :
une nouvelle évaluation incorpore l’état, son registre et sa signature, la section
de tour, le lecteur, la réalisation, la coordinatisation de comparaison et la
provenance. Les cas antérieurs conservent leur contenu et leur date
d’arrêté. Ainsi, la croissance de l’inventaire rend visible une plus grande partie
du domaine sans altérer rétrospectivement ce que chaque preuve et
chaque évaluation avaient établi.

